\section{Conclusion and Future Work}\label{sec:conclusion}
In conclusion, this work presents a comprehensive study on using deep reinforcement learning to develop versatile, robust, and dynamic locomotion controllers for bipedal robots. 
This work introduces a dual-history approach that integrates the robot's input/output (I/O) history into RL-based controllers and highlights its significance. 
We demonstrate the adaptivity and robustness of the proposed RL-based controller, particularly underscoring how a well-designed long I/O history encoder can adapt to both time-invariant dynamics changes and time-variant events such as varying contacts. 
Additionally, task randomization, which encourages the robot to explore a broader range of scenarios by accomplishing different tasks, significantly enhances robustness, complementing the robustness achieved through traditional dynamics randomization.
The proposed method is validated thoroughly on the bipedal robot Cassie, successfully realizing versatile and robust walking, running, and jumping skills in the real world. 
These experiments include several state-of-the-art results, including walking control with consistent performance over a long timespan (459 days), versatile running capabilities demonstrated in a 400-meter dash and across challenging terrains, and a large repertoire of jumping tasks including the furthest 1.4-meter forward jump and 0.44-meter high jump.


In the future, we hope this method can be extended to humanoid robots that can also leverage upper-body motions for agility and stability. 
Additionally, integrating depth vision directly into the locomotion controller can be realized in a straightforward way where an additional depth encoder can work alongside the robot's I/O history encoder in our proposed control architecture.  
We encourage the readers to also try this setup for visual control on their bipedal robots. 
With the advancements made in this work, we think a large portion of practical problems in realizing effective locomotion control for human-sized bipedal robots can be addressed. 
For the bipedal robotics community, an exciting direction would be the combination of bipedal locomotion and bimanual manipulation to tackle long-horizon loco-manipulation tasks, opening new possibilities in the field.